\chapter{Discussion, Limitations, and Future Work}
\label{ch:discussion}

\section{Discussion}
\label{sec:disc-discussion}

Read across Chapters~5 through 9, AI Compass's central architectural claim is not that it avoids generative AI, but that it confines it to a small number of well-defined, independently verifiable roles and keeps it out of the two subsystems where an unverifiable model output would be most costly: ranking and retrieval. Content enrichment (Chapter~5) uses exactly one structured LLM call per item, with defensive validation that can only narrow or reject a model's output, never invent a category or entity outside the controlled vocabulary. RAG answer generation (Chapter~7) and long-video chapter summarization (Chapter~8) are the only other places a language model is invoked, and both are wrapped in mechanical, server-side checks (citation-marker resolution in the first case, map-reduce chunk boundaries snapped to transcript segments in the second) that constrain what a generative step is allowed to assert rather than trusting it outright. The ranking core (Chapter~6) and the retrieval mechanics underneath both the recommender and the RAG assistant (cosine similarity over stored embeddings, Chapters~6 and~7) are, by contrast, fully deterministic: no step in either path issues a call to a language model, a property confirmed directly by code inspection in Chapter~10 rather than merely asserted.

This separation has two direct, checkable consequences that recur through the thesis. First, the personalized ranking is fully explainable: every relevance score is a closed-form function of five documented signals (Chapter~6), and every explanation shown to a user is produced by template selection over the same signals, not by a second, independent generative pass that could drift from what was actually scored. Second, RAG grounding is enforced structurally rather than aspirationally: the system does not merely instruct the model to cite its sources and hope, it parses the model's citation markers after the fact and marks an answer ungrounded whenever every marker in it fails to resolve against the retrieved context (Chapter~7). Neither property required trusting a language model to behave; both are consequences of where, in the pipeline, the deterministic and the generative components were drawn to meet.

The same discipline extends to how the system handles its own retrieval infrastructure. Access control on retrieved RAG passages is applied as a mechanical filter over source and category visibility before any passage reaches the generation step (Chapter~7), not as an instruction folded into the prompt and hoped for; a passage the requesting user is not entitled to see is removed from the candidate set before the model ever sees it, so a leak would require a bug in the filter, not a lapse in the model's judgment. Chapter~10's structural verification of embedding-write idempotency, a database-level uniqueness constraint holding even under a documented past race condition, is another instance of the same pattern: correctness properties that matter are pushed down to a layer that can guarantee them mechanically, rather than left to the discipline of whichever code path happens to run first. Taken together, these choices amount to a consistent engineering stance rather than a collection of unrelated decisions: wherever a property needs to be trusted, the system is built so that property is instead checkable.

\section{Limitations}
\label{sec:disc-limitations}

\subsection{Data and ingestion layer}
\label{sec:disc-lim-data}

Two structural gaps constrain how AI Compass will behave as its corpus grows. First, the general-purpose \texttt{embeddings} table carries no approximate-nearest-neighbor index (only the RAG-specific \texttt{rag\_chunks} table has one), so every clustering neighbor lookup, ranking similarity search, and semantic search against \texttt{embeddings} is currently an exact scan of the whole table rather than an indexed approximate search. Second, the pipeline's embedding phase only scans the newest 1000 items per content type on each run, with no offset-based pagination; once the corpus exceeds that figure, an item that failed to embed while it was new is never revisited, a real scale ceiling rather than a hypothetical one. Third, ingestion writes use \texttt{ON CONFLICT DO NOTHING} at the database level, which makes ingestion cheaply idempotent but means a corrected or expanded version of an already-seen article or video is never re-ingested; the stale first version is what the system keeps indefinitely.

\subsection{Ranking limitations already disclosed in Chapter~6}
\label{sec:disc-lim-ranking}

Two gaps in the personalized ranking formula are documented in detail in Section~\ref{sec:rec-limitations} and are restated here only for completeness. The exploration slice's weighted-random draw uses Python's unseeded \texttt{random.random()}, so which leftover candidates fill the exploration slots is not reproducible across repeated ranking runs over an identical candidate pool for the same user. Separately, \texttt{UserAffinity}'s entity dimension is computed and stored nightly but is never read by the scoring formula itself, which consumes only the topic and source dimensions, a real, unaddressed gap between what is computed and what personalization actually uses.

\subsection{Verification and operational coverage}
\label{sec:disc-lim-ops}

Automated test coverage is thin precisely where the logic is hardest to get right: there is no test suite exercising \texttt{RankingService}, the RAG retrieval or generation path, or the near-duplicate clustering logic, so correctness claims about these subsystems rest on direct code inspection and manual verification rather than on a regression suite that would catch a future change breaking them silently. Separately, no Celery task in the pipeline declares a retry policy (no \texttt{autoretry\_for}, backoff, or maximum-retry setting anywhere), so a task that raises an exception is simply lost until the next scheduled run happens to reprocess the same work idempotently, rather than being recovered automatically; the speech-to-text dispatch path shows a concrete instance of this gap, since a job flipped from \texttt{queued} to \texttt{running} before dispatch has no mechanism to return to \texttt{queued} if the worker that claimed it dies, and can therefore remain stuck indefinitely. Relatedly, no monitoring or alerting layer observes any of this: if a scheduled beat process stops, if a worker dies, or if the LLM provider begins rejecting requests, nothing surfaces the failure, since the Celery task wrapper around the pipeline returns normally regardless of whether the underlying run succeeded. Finally, the evaluation dataset's scale is itself a limitation of what could be measured, not of what was built; Chapter~10 details why no meaningful NDCG@10 or MAP number could yet be produced.

\section{Future Work}
\label{sec:disc-future}

The items below are disclosed scope decisions rather than gaps: each was left for later because the constraint that motivates it is already documented elsewhere in this thesis, and each has a concrete, bounded path forward.

The most immediate is re-running the offline ranking evaluation once AI Compass has accumulated a period of sustained real usage, this time with a held-out window matched to the ranking system's actual three-hour refresh cadence rather than to an arbitrary multi-day span, exactly as Chapter~10 recommends. Adding an approximate-nearest-neighbor index to the \texttt{embeddings} table, mirroring the HNSW index \texttt{rag\_chunks} already has, would remove the exact-scan ceiling described in Section~\ref{sec:disc-lim-data} with a comparatively small, well-scoped migration. Replacing the current \texttt{ON CONFLICT DO NOTHING} ingestion strategy with a content-hash comparison (re-ingesting, and invalidating downstream enrichment and embeddings for, an item whose hash has changed) would close the lost-update gap in the same subsection without sacrificing ingestion idempotency for unchanged content. Incorporating the currently-computed-but-unused entity-affinity signal into the personalized ranking formula (Chapter~6, Section~\ref{sec:rec-limitations}) is a bounded change to an already-identified gap rather than new design work. Collaborative filtering remains deferred until the active user base is large enough for cross-user co-engagement signal to be statistically meaningful, the same cold-start reasoning already given in Chapters~2 and~6. Per-persona summaries and region- or language-aware filtering (Chapter~6) remain future scope for the same reasons stated there: the first would grow the enrichment budget from per-item to per-user, and the second is not yet justified by a content pool or user base diverse enough to make such filters meaningful.

The operational gaps identified in Section~\ref{sec:disc-lim-ops} point to similarly bounded fixes rather than open-ended redesign. Adding a modest Celery retry policy (bounded backoff and a maximum retry count on the tasks that call external services) would recover the class of failures that currently require waiting for the next scheduled run, and a small reaper task that resets speech-to-text jobs stuck in \texttt{running} back to \texttt{queued} after a fixed timeout would close the specific stuck-job path described above. A basic health and alerting layer over the scheduled beat process, the worker pool, and the LLM provider's response codes would turn currently-silent failures into visible ones, which the system today has no mechanism for at all. None of these require touching \texttt{RankingService}, the RAG pipeline, or any scoring formula described earlier in this thesis; they are operational hardening around a design whose core logic this thesis has argued is already sound.
